%%%%%%%%%%%%%%%%%%%%%%%%%%%%%%%%%%%%%%%%%%%%%%%%%%%%%%%%%
%    Formato de memoria de Trabajo Final de Grado (TFG) para la Facultad de Ingeniería de la Universidad Nacional de Asunción.

%   Version 2.0.1 por Msc. Ing. Sergio Ramón Toledo Gallardo del 10 de junio de 2015                           
%- La compilación debe realizarse utilizando el comando pdflatex o pdflatexmk (para Mac)%
% Los programas necesarios son Miktex 2.9 (distribución latex) y algún procesador como el texmaks o texniccenter
%- Para texnicCenter Usar makeindex y habilitar usar bibtex
%-  Al imprimir el PDF se debe seleccionar tamaño real y Elegir origen por tamaño de pagina PDF para respetar los margenes requeridos   
%- Las figuras deben cargarse en la carpeta "imagenes" en formato pdf para optimizar la velocidad de compilación, sin embargo también se permiten formatos jpg y eps.
%- Las referencia bibliográficas se cargan en el archivo    fiuna.bib                                                       %%%%%%%%%%%%%%%%%%%%%%%%%%%%%%%%%%%%%%%%%%%%%%%%%%%%%
\documentclass[oneside, a4paper, 12pt]{fiuna}
\usepackage{fiuna}

%%%%%%%%%%%%%%%%%%%%%%%%%%%%%%%%%%%%%%%%%%%%%%%%%%%%%%%%%%
%               Inicio del documento                     %
%%%%%%%%%%%%%%%%%%%%%%%%%%%%%%%%%%%%%%%%%%%%%%%%%%%%%%%%%% 
\begin{document}

%%%Introduzca el nombre de los autores
\autor[Nombre Corto]{Nombre completo del autor}
\autor[Sergio Toledo]{Sergio Ramon Toledo Gallardo}
%\autor{Autor 2}
%%%Introduzca el título del TFG
\titulo{Titulo de TFG}
%%%introduce la carrera
\carrera{Ingenier\'ia en Algo}
%%%Nombre del Asesor del TFG %%%%%
\asesor{Ing. Asesor}
%\asesor{Ing. Asesor 2}
%\asesor{Ing. Asesor 3}
%%%%%%Ciudad donde se presenta%%%%%%%%%%%
\ciudad{San Lorenzo}
%%%%%% A quien se dedique %%%%%%%%%%%%%%
\dedicatoria{A quien se dedique}
\pagenumbering{roman}

\preliminares

\pagenumbering{arabic}
\tableofcontents
\listoffigures
\listoftables
%% En caso de no tener anexos comentar la siguiente linea
\listofappendices

\formatoTitulos

%%% Inician los capítulos de la memoria de TFG%%%%
%cada capitulo se introduce con \input
%El primer capitulo siempre es la introducción al TFG
\input{capitulo01}
%Luego inican los demas capítulos
\input{capitulo02}    
%%%bibliografía%%%%%%%%%
%Cargar los libros en formato bibtex en el archivo fiuna.bib
\formatoIndice
\bibliographystyle{fiuna}
%\nocite{*}
\bibliography{fiuna}   

%%%anexos%%%%%%
%En caso de no tenerlos, comentar la siguiente línea
\input{anexo}
\end{document}